\documentclass[compacto,10pt]{pubciencias-notas}

% -- Paquetes adicionales
\usepackage{comandosEPN}

% -- Datos del libro
\autor{Departamento de Formación Básica}
\materia{Álgebra Lineal}
\nota{Semana 12 Clase 01}
\fecha{Semestre 2024-B}

\logouno[1.1cm]{Logo_EPN.eps}
\logodos[1.7cm]{Logo_DFB.eps}

\definecolor{colordef}{cmyk}{0.81,0.62,0.00,0.22}

\begin{document}

\encabezado
\vspace*{-8mm}

\setcounter{section}{4}
\section{Espacios con producto interno}
\setcounter{subsection}{2}
\subsection{Proyección ortogonal}
\setcounter{defi}{8}
\begin{defi}[Proyección ortogonal]
Sean $(E,+,\cdot,\K)$ un espacio vectorial con producto interno, 
$W$ un subespacio vectorial de $E$ y  $\{u_1, u_2, \ldots, u_k \}$ una base ortogonal de $W$. Para $v \in E$, la proyección ortogonal de $v$ sobre $W$, denotado por $\proy_W (v)$, se define por:

\[   \proy_W(v) = \dfrac{\left<v,u_1 \right>}{\left<u_1, u_1 \right\rangle }u_1 +
   \dfrac{\left<v,u_2 \right\rangle}{\left<u_2, u_2 \right\rangle}u_2 +
   \cdots +
   \dfrac{\left<v,u_k \right\rangle}{\left<u_k, u_k \right\rangle}u_k,
\]
donde $\proy_W (v) \in W$.
\end{defi}
\begin{ejem}\hspace{0pt}\label{ProyeccionOrtogonal} En el espacio vectorial $\R^{3}$ considere a $B=\{u_1, u_2\}$ una base ortonormal de un subespacio $W$ de $\R^3$, donde
\[
    u_1 = \l(\dfrac{\sqrt{2}}{2},\dfrac{\sqrt{2}}{2},0\r)
    \texty
    u_2 = \l(- \dfrac{\sqrt{6}}{6},\dfrac{\sqrt{6}}{6}, \dfrac{\sqrt{6}}{3}\r)
\]
y sea $v = (1,1,1)$, entonces
\begin{align*}
\proy_W(v)
&= \l<v, u_1\r> u_1 + \l<v, u_2\r> u_2\\
&= \l<(1, 1, 1), \l(\dfrac{\sqrt{2}}{2},\dfrac{\sqrt{2}}{2},0\r)\r> \l(\dfrac{\sqrt{2}}{2},\dfrac{\sqrt{2}}{2},0\r) +\\
&\phantom{= } \l<(1, 1, 1), 
\l(- \dfrac{\sqrt{6}}{6},\dfrac{\sqrt{6}}{6}, \dfrac{\sqrt{6}}{3}\r)\r> \l(- \dfrac{\sqrt{6}}{6},\dfrac{\sqrt{6}}{6}, \dfrac{\sqrt{6}}{3}\r)\\
&= \sqrt{2}\l(\dfrac{\sqrt{2}}{2},\dfrac{\sqrt{2}}{2},0\r) + \dfrac{\sqrt{6}}{3} 
\l(- \dfrac{\sqrt{6}}{6},\dfrac{\sqrt{6}}{6}, \dfrac{\sqrt{6}}{3}\r)\\
&= \l(1, 1, 0\r) + \l(- \dfrac{1}{3}, \dfrac{1}{3}, \dfrac{2}{3}\r)\\
&= \l(\dfrac{2}{3}, \dfrac{4}{3}, \dfrac{2}{3}\r)
\end{align*}
donde $\proy_W(v) \in W = \{(x_{1}, x_{2}, x_{3})\in\R^{3}\,:\,x_{1} - x_{2} + x_{3} = 0\}$.
\end{ejem}
\begin{teo}[Teorema de la proyección]\label{TeoProyec}
Sean $(E,+,\cdot,\K)$ un espacio vectorial de dimensión finita provisto de producto interno, $W$ un subespacio vectorial de $E$ y $v \in E$. Se tiene que
\[
    v = \proy_W( v) + \proy_{W^\perp}( v).
\]
\end{teo}

\begin{ejem}\hspace{0pt} Considere el ejemplo \ref{ProyeccionOrtogonal}, se puede encontrar directamente la proyección ortogonal de $v$ sobre $W^{\perp}$, aplicando el teorema \ref{TeoProyec}:
\begin{align*}
\proy_{W^\perp}(v) 
    &=  v -\proy_{W}(v)\\
    &= (1,1,1) - \l(\dfrac{2}{3}, \dfrac{4}{3}, \dfrac{2}{3}\r)\\
    &= \l(\dfrac{1}{3}, - \dfrac{1}{3}, \dfrac{1}{3}\r)
\end{align*}
sin necesidad de determinar previamente que:
\[
W^{\perp} = \{(x_{1}, x_{2}, x_{3})\in\R^{3}\,:\,x_{1} + x_{2} = 0, x_{2} + x_{3} = 0\}.
\]
\end{ejem}
\begin{teo}
Sean $(E,+,\cdot,\K)$ un espacio vectorial de dimensión finita provisto de producto interno, $W$ un subespacio vectorial de $E$ y $v \in E$. Se tiene que el vector en $W$ más cercano a $v$ es $\proy_W(v)$, es decir,
\[
\|v-w\| \quad \text{ es mínima cuando } \quad w = \proy_W(v).
\]
Equivale a afirmar que, para todo $w\in W$ se verifica la desigualdad:
\[
    \|v - \proy_{W}(v)\| \leq \|v - w\|,
\]
y se cumple la igualdad solo en el caso en que $w = \proy_{W}(v)$.
\end{teo}

\begin{ejem}\hspace{0pt} Considere el ejemplo \ref{ProyeccionOrtogonal}, se tiene que tomando $w = \proy_{W}(v)$ el valor de $\|v - w\| = \|\proy_{W^{\perp}}(v)\| = \dfrac{\sqrt{3}}{3}$ es la distancia mínima de $v = (1,1,1)$ a $W = \{(x_{1}, x_{2}, x_{3})\in\R^{3}\,:\,x_{1} - x_{2} + x_{3} = 0\}$.
\end{ejem}

\newpage
\section{Aplicaciones lineales}

\subsection{Definición de Aplicación lineal}
\setcounter{defi}{0}
\begin{defi}[Aplicación lineal]
    Sean $(E,+_1,\cdot_1,\K)$ y $(F,+_2,\cdot_2,\K)$ espacios vectoriales. A una función
    
    \noindent $\func{T}{E}{F}$ se le llama aplicación lineal (transformación lineal) si para todo $\alpha \in \K$, y todo $u,v \in E$ se cumple:
    \begin{enumerate}
    \item 
        $T(u +_1 v) = T(u) +_2 T(v)\quad$ y
    \item 
        $T(\alpha \cdot_1 v ) = \alpha \cdot_2 T(v)$.
    \end{enumerate}
\end{defi}

\begin{obs}
En adelante, consideraremos $(E,+,\cdot,\K)$ y $(F,+,\cdot,\K)$ espacios vectoriales, ya que, si bien se están notando de la misma manera las operaciones suma y producto de los dos espacios vectoriales, estamos claros que no tiene por qué ser así. Con esto conseguimos simplificar la escritura de la definición de aplicación lineal.
\end{obs}

\begin{advertencia}
    Se denotará
    \begin{enumerate}
    \item 
        $T(u + v) = T(u) + T(v)\quad$ y
    \item 
        $T(\alpha v ) = \alpha T(v)$,
    \end{enumerate}
    para $\alpha \in \K$ y $u,v \in E$.
\end{advertencia}

\begin{obs}
    Al conjunto de todas las aplicaciones lineales del espacio vectorial $E$ en el espacio vectorial $F$, lo designaremos por $\mathscr{L}(E,F)$.
\end{obs}

\begin{ejem}\hspace{0pt} \label{ejem1}
Las siguientes funciones, son ejemplos de transformaciones lineales.
\begin{enumerate}
\item
    Aplicación lineal nula.
    \[
        \funcion{T}{E}{E}{v}{0_{E}}
    \]
\item
    Aplicación lineal identidad.
    \[
        \funcion{T}{E}{E}{v}{v}
    \]
    \item
    \[
        \funcion{T}{\R^3}{\R^2}{(x,y,z)}{(x-2z,y-z)}
    \]
      \item
    \[
        \funcion{T}{\R^3}{\R^3}{(x,y,z)}{(x + y,y + z,x)}
    \]
  
    
   
    
  
\end{enumerate}
\end{ejem}
\begin{ejem}
En la aplicaci\'on lineal
\[
        \funcion{T}{\R^3}{\R^2}{(x,y,z)}{(x-2z,y-z)}
\]
se tiene que
\[
    T(0,0,0)=(0,0) \texty T(-1,2,3)=(-7,-1).
\]
\end{ejem}

\begin{ejer}\label{Ejercicio1}
Dada la aplicaci\'on lineal 
    \[
        \funcion{T}{\R^3}{\R^3}{(a,b,c)}{(a,a-b,c)}
    \]
Hallar  $T(0,0,0)$, $T(1,-2,3)$ y $T(0,-6,1)$.
\end{ejer}


\begin{ejem}
Demostremos que la función definida en el Ejercicio \ref{Ejercicio1} es una aplicación lineal.
\[
    \funcion{T}{\R^3}{\R^3}{(a,b,c)}{(a,a-b,c)}
\]
\begin{enumerate}
\item
    Para todo $u, v \in \R^{3}$, vamos a demostrar que $T(u+v) = T(u) + T(v)$.

    \noindent Sean $u = (a, b, c), v = (d, e, f) \in \R^{3}$, tenemos
    \begin{align*}
        T(u) & = (a,a-b,c) \\
        T(v) & = (d,d-e,f) \text{, luego} \\
        T(u) + T(v) & = (a,a-b,c) + (d,d-e,f).
    \end{align*}
    Por otra parte, $u+v = (a+d, b+e, c+f)$, se sigue que
    \begin{align*}
        T(u+v) & = (a + d,(a+d) - (b+e), c+f) \\
        & = (a,a-b, c) + (d,d-e,f).
    \end{align*}
    Por lo tanto:
    \[
        T(u+v) = T(u) + T(v)
    \]
    
\item
    Para todo $\alpha \in \R$ y para todo $u \in \R^{3}$, vamos a demostrar que $T(\alpha u) = \alpha T(u)$.
    
    Sean $\alpha \in \R$ y $u = (a, b, c) \in \R^{3}$. 
    Tenemos que $\alpha u = (\alpha a, \alpha b, \alpha c)$, luego
    \begin{align*}
    T(\alpha u) & = (\alpha a,(\alpha a - \alpha b) , \alpha c) \\
    & = \alpha (a, a-b, c)
    \end{align*}
    Por otra parte
    \[
        \alpha T(u) = \alpha (a ,a-b,c)
    \]
    Por lo tanto:
    \[
        T(\alpha u) = \alpha T(u)
    \]
    De $1.$ y $2.$ se sigue que $T$ es una aplicación lineal.
\end{enumerate}
\end{ejem}

\begin{ejer} 
Demuestre que la siguiente funci\'on es una aplicaci\'on lineal:
   \[
        \funcion{T}{\R^3}{\R^2}{(x,y,z)}{(x-2z,y-z)}
    \]
\end{ejer}
\begin{teo}
     Sea $T \in \mathscr{L}(E,F)$. Para todo $u,v, v_1, v_2, \ldots, v_n \in E$ y $\alpha_1, \alpha_2, \ldots, \alpha_n \in \K$, se tiene que:
     \begin{enumerate}
         \item $T(0_E) = 0_F$.
         \item $T(u-v) = T(u) - T(v)$.
         \item $T\left(\displaystyle \sum_{k=1}^n \alpha_k v_k \right)= 
         \displaystyle \sum_{k=1}^n \alpha_k T\left(v_k \right)$.
     \end{enumerate}
\end{teo}

\begin{teo}
    $(\mathscr{L}(E,F), +, \cdot, \K)$ es un espacio vectorial. Siendo $+$ la suma usual entre funciones y $\cdot$ la multiplicación de un elemento del campo por una función.
\end{teo}
\subsection{Propiedades de las aplicaciones lineales}

\begin{prop}
Sean $T_1,T_2 \in \mathscr{L}(E,F)$ y $B= \{v_1, v_2, \ldots, v_n\}$ una base para $E$. Si
\[
    T_1(v_i) = T_2(v_i)
\]
para todo $i \in \{1, 2, \ldots, n\}$, entonces $T_1(v)=T_2(v)$, para todo $v\in E$, es decir, 
\[
    T_1 = T_2.
\]
\end{prop}

\begin{prop}
Si $\{v_1, v_2, \ldots, v_n\}$ es una base del espacio vectorial $E$ y 

\noindent \{ $w_1, w_2, \ldots, w_n \} \subset F$, entonces existe una única transformación lineal

\noindent $T \in \mathscr{L}(E,F)$ tal que
\[
    T(v_i) = w_i
\]
para todo $i\in \{1, 2, \ldots, n\}$.
\end{prop}


\end{document}
